\section{Source audit, proposed qualifications, and verification}
\label{sec:audit}

\subsection{A fixed source record}

The main source baseline is the ProveIt commit whose full identifier is
\begin{center}
\small\src{921b58f1099568bfefa28a556470e4b29e94eb3b}.
\end{center}
It contains the manuscript directory supplied for this task and four
incoming archives dated 10--11 October 2026: the bilinear
Mellin--Hurwitz, cubic harmonic mixed-order, exact resonant, and
independent-order reports \cite{RepoMain,RepoBilinear,RepoCubic,
RepoExactResonant,RepoIndependent}. The earlier mixed spectral report,
which states several questions answered here, was inspected at the
historical revision
\begin{center}
\small\src{5a790187c8e186e41e2b990b4941cb7a1a3c7b6b}.
\end{center}
Its archived source is cited separately as \cite{RepoMixed}; earlier
exact identities used in the shift calculation are \cite{RepoExact}.
This distinction matters because an incoming archive can be integrated
and removed from a later incoming directory without invalidating its
historical role.

The audit covered the definitions, proofs, status statements, and
questions directly used in this continuation, together with the current
Gaussian conjecture labels. It did not certify every line of the entire
multi-volume manuscript. No new false theorem was identified in the
selected source material. The qualifications below state how the new
results must be recorded; they are not allegations that the source
already contains each invalid inference.

\subsection{Status changes justified by the proofs}

\paragraph{The shift-boundary question is answered.}
The mixed spectral report asks for the actual germ at $a=-r$ and the
exception at nonpositive spectral orders. Theorems~\ref{hsh:germ},
\ref{hsh:boundary}, \ref{hsh:radius}, and \ref{hsh:rank} answer it.
The existing half-plane theorem remains correct, but a status summary
can now give the exact first singularity and exact Taylor radius.
Transport itself should be attributed to the broader relative ordered
Hurwitz character: the exact identification is
\eqref{hsh:relativecharacter}. Treating the present family as unrelated
to that source would duplicate an established construction.

\paragraph{The weight-six and weight-eight separation question is answered
within its specified family.}
Equations~\eqref{bell:sixrref} and \eqref{bell:eightrref} give the
requested rational elimination. Theorem~\ref{bell:rank} supplies the
general rank, while Theorem~\ref{bell:image} determines the full
summation kernel. The individual nonlinear coordinate moments at
weights six and eight lie outside this formal mixed-Bell span.
This does not rule out evaluating them using another parameter family,
an argument derivative, or an independent functional identity.

\paragraph{The compact harmonic-power target has a constructive solution.}
Theorem~\ref{hpnew:generator} and
Corollary~\ref{hpnew:threerows} answer the complete minus-two sequence
requested in Question~10 of \cite{RepoExactResonant}; they also give
every further centered even moment. Theorem~\ref{hpnew:uniformspan}
adds the stronger conclusion that one finite Euler-sum span works for
all moment orders at a fixed power. A closed bivariate expression and
minimal arithmetic coordinates remain open.
Corollary~\ref{hpnew:oddorders} supplies a sufficient generalized-harmonic
parity class toward Question~11 and a finite formula for every single
positive odd harmonic order; it does not finish the requested necessity
classification.

\paragraph{Nested collision constants are now explicit in the polygamma
sector.}
The exact formula of Theorem~\ref{hc:main} supplies a jointly
holomorphic remainder for any number of factors and any nonnegative
argument-derivative orders. Theorem~\ref{hc:hierarchy} answers the
named nested triple paths and all their monomial-rate analogues.
The extension to positive Stieltjes indices and a theorem about
composition of renormalized merging operations are still separate
questions. They should not be marked solved by the present result.

\subsection{Qualifications that prevent incorrect extensions}

\paragraph{Principal continuation is a statement about a branch.}
The principal family is holomorphic on the slit plane in
\eqref{hsh:slit}. Earlier removable shifts need not remain removable
on every sheet. A loop around $a=-1$ changes $D_1$ by
\[
 \bigl(e^{-2\pi\ii s}-1\bigr)\frac{(a+1)^{-s}}a.
\]
The added term generally has a pole at $a=0$. A claim of unrestricted
removability on the universal continuation would therefore need
correction. The source's valid half-plane result makes no such claim.

\paragraph{Polynomial values do not imply polynomial spectral derivatives.}
The finite Stirling formula at $s=-m$ has rank $m+1$ in depth.
Differentiating the analytic family introduces
$z^m(-\log z)^k$ at the first boundary. For every $k\ge1$ it has
finite shift Taylor radius and supplies independent depth functions.
It cannot be obtained by differentiating a formula whose label $m$
has only been specified at integers. Similarly, the first spectral
derivative of the raw harmonic-power family retains a global Mellin
remainder that is absent from its regular value.

\paragraph{A removable restriction need not be a jointly removable germ.}
The auxiliary $G_m$ in Section~\ref{hpnew:section} obeys $G_m(0)=0$,
whereas separately specializing the uncentered family at exponent zero
gives $\zeta(-m)$. The explicit term in
Remark~\ref{hpnew:orderwarning} proves the discrepancy. Its polar
hyperplane is not removable merely because its numerator vanishes at
one point. In contrast, the centered family has no polar hyperplane
through $(-2r,0)$, so the claimed coefficient extraction is valid.

\paragraph{A collision constant depends on the full primitive.}
The terms of order $p\ge2$ in \eqref{hc:counter} are rational
primitive terms, in addition to logarithms. They are essential:
the derivative example \eqref{hc:middle-derivative} would be wrong
without them. Even a curved binary gap $\epsilon+b\epsilon^2$
changes the raw constant by $b$, as
\eqref{hc:binary-curved} shows. Normalizing only a leading pole or
only a fixed-rate approximation does not determine the required
constant along a nested path.

\paragraph{Formal independence and arithmetic independence are distinct.}
The Bell rank treats the $L_k$ as independent indeterminates, and the
shift rank treats analytic functions of $a$ over the constants.
The Euler sums in $\mathcal E_{p-1}$ form a proved spanning list.
None of these assertions gives a lower bound for the arithmetic
dimension of their numerical values. The exact matrices are suitable
certificates for the specified algebraic problem and no larger one.

\paragraph{The retained Gaussian conjectures remain conjectures.}
The pinned manuscript labels the current $S_6$ candidate
\src{cycloquot:conj:S6} in
\src{chapters/04-cyclotomic-quotients.tex}, and the revised $S_8$
candidate \src{s8new:conj:S8} in
\src{chapters/04-S8-candidate.tex} \cite{RepoMain}.
This article proves neither numerical period identity. The rejected
older $S_8$ vector is already distinguished in the current source.
No new integer-relation fit or conjectural coefficient vector is
offered here. The remaining task is a new exact functional relation
with its analytic endpoint evaluation.

\subsection{Reproducible exact checks and independent diagnostics}

Every primary analytic result has a proof in the preceding sections.
The following additional checks are included as readable Python source
and JSON results. They require only Python, SymPy, and mpmath.
The driver \src{code/run_all.py} runs the checks in a fixed order.

\begin{longtable}{@{}p{0.24\textwidth}p{0.69\textwidth}@{}}
\toprule
Family & Recorded verification\\
\midrule
\endhead
Shift germs & 85 exact core checks and 50 numerical comparisons:
Newton versus defining sums, finite transport, Cauchy-extracted
boundary constants, spectral jets, and a doubled cutoff.
The primitive supplement adds 12 exact derivative and partial-fraction
checks.\\[5pt]
Mixed Bell moments & Independent coefficient constructions and ranks
for weights $2$ through $18$; exact elimination and annihilators at
weights $6,8$; the quartic zero row; and a deliberately corrupted
coefficient that is correctly rejected. Nineteen mixed rows at five
real or complex parameter choices are checked numerically, plus the
independent harmonic tail-square identity.\\[5pt]
Harmonic powers & Seven exact rational generator residuals;
26 Bernoulli identities through moment order $r=12$;
15 centered values at powers $1$--$5$ and moment orders $0$--$2$;
and six full exponent-generator comparisons at two exponent values.
The centered computations use two truncation lengths.\\[5pt]
Generalized harmonics & Nine exact centered Faulhaber identities and
six independent values for harmonic orders $q=3,5$ at moment orders
$r=0,1,2$, also using two subtraction truncations.\\[5pt]
Collisions & 132 exact oriented hierarchy assertions for
$1\le q<p\le12$, exact quartic coefficients and rate differences,
and the derivative constant. Independent local-subtraction quadrature
is used for 12 digamma configurations and two polygamma configurations.
The corrected remainder tends to the merged value as predicted.\\
\bottomrule
\end{longtable}

The largest scaled residual in the core shift comparisons is about
$1.76\cdot10^{-23}$ and is controlled by the chosen Newton cutoff.
The Bell mixed-row comparisons are within
$2.70\cdot10^{-65}$ in the recorded scaled metric; the tail-square
comparison has absolute residual below $8.30\cdot10^{-75}$.
For the 15 centered harmonic-power values, the largest absolute
discrepancy is below $6.49\cdot10^{-45}$; the independent truncations
differ by at most $1.29\cdot10^{-32}$. These are recorded diagnostic
values, not rigorous error estimates. The precision, truncation
parameters, and individual residuals are retained in the result files.

The collision diagnostics test the exact correction
$\mathcal C-\mathcal A$ against its proved holomorphic limit.
They do not assert that this difference equals its limit at a nonzero
shift: the remaining Taylor terms are expected. This makes their
numerical interpretation different from the residual of an exact
finite-parameter Bell identity.

\subsection{Integration order}

The source sections can be integrated independently, retaining their
label prefixes \src{hsh:}, \src{bell:}, \src{hpnew:}, and \src{hc:}.
First reconcile the relative-character identification, then add the
shift singularities and primitives; the Bell and harmonic-power
sections can follow their existing antecedents; finally insert the
collision theorem before the nested-rate examples. The plain-text
integration note records exact source dependencies and suggested
status replacements. The complete standalone TeX and its compiled PDF
are included, alongside the modular source and verification programs.
